\documentclass[12pt]{article}

\usepackage[T1]{fontenc}
\usepackage[utf8]{inputenc}
\usepackage[english]{babel}

\usepackage{lmodern}
\usepackage{geometry}
\usepackage{graphicx}
\usepackage{fancyhdr}
\usepackage{amsmath}
\usepackage{tocloft}
\usepackage{titlesec}
\usepackage{hyperref}
\usepackage{float}
\usepackage{listings}
\usepackage{xcolor}

\hypersetup{
    colorlinks=true,
    linkcolor=black,
    urlcolor=black
}

\renewcommand\thesection{\Roman{section}}
\renewcommand\thesubsection{\thesection.\roman{subsection}}

\geometry{
    a4paper, top=1in, bottom=1in, left=1in, right=1in,
    headheight=15pt, footskip=15pt,
}

\lstset{
    language=Matlab, basicstyle=\ttfamily\small,
    keywordstyle=\color{blue}, commentstyle=\color{gray}, stringstyle=\color{red},
    numbers=left, numberstyle=\tiny, stepnumber=1, numbersep=5pt,
    backgroundcolor=\color{white}, showspaces=false, showstringspaces=false,
    showtabs=false, frame=single, tabsize=4, captionpos=b, breaklines=true
}

\linespread{1.06}
\setlength{\parindent}{0pt}
\setlength{\parskip}{0.8em}

\pagestyle{fancy}
\fancyhf{}
\lhead{DVD-Assignment 00}
\cfoot{\thepage}
\renewcommand{\headrulewidth}{0.4pt}
\renewcommand{\footrulewidth}{0pt}

\titlespacing*{\section}{0pt}{8pt}{6pt}
\titlespacing*{\subsection}{0pt}{6pt}{4pt}

\newcommand{\includelogo}{
    \IfFileExists{IIIT_Hyderabad_Logo.jpg}{%
        \includegraphics[width=.8\textwidth]{IIIT_Hyderabad_Logo.jpg}%
    }{%
        \fbox{\parbox[c][4cm][c]{0.8\textwidth}{\centering\Large Image placeholder\\[0.5em]\small IIIT Hyderabad Logo}}%
    }%
}

\newcommand{\includeimage}[3][0.75\linewidth]{%
    \IfFileExists{#2}{%
        \includegraphics[width=#1]{#2}%
    }{%
        \fbox{%
            \parbox[c][4.5cm][c]{#1}{%
                \centering
                \Large Image Placeholder\\[0.5em]
                \small #3 (\texttt{#2})
            }%
        }%
    }%
}

\begin{document}

\begin{titlepage}
    \centering
    \vspace*{12mm}
    \includelogo
    \vspace{10mm}
    {\large\scshape International Institute of Information Technology\\ Hyderabad \par}
    \vspace{6mm}
    {\large \bfseries Design VLSI Design  \par}
    \vspace{8pt}
    {\large \bfseries Assignment 0 \par}
    \vspace{6pt}
    {\large \textit{Course Code: M26EC2.408} \par}
    \vfill
    {\Large\scshape \bfseries Ravi Kumar Sahu \& Diddigi Srivatsa  \par}
    \vspace{4pt}
    {\large Roll No: 2024102024  \& 2025122009\par}
    \vspace{10mm}
    \thispagestyle{empty}
\end{titlepage}

\clearpage
\pagenumbering{roman}
\setcounter{page}{1}
\thispagestyle{plain}
\tableofcontents
\clearpage

\pagenumbering{arabic}
\setcounter{page}{1}
\pagestyle{fancy}
\fontsize{11}{13}\selectfont

\section{Part A — RTL Design and Verification}

\subsection{Design Overview}
The design implements a 4-bit unsigned array multiplier generating an 8-bit product using registered inputs ($a$, $b$) and a registered output ($p$). The datapath is built strictly from an array of 16 AND gates generating partial products and 12 structural instances of a 1-bit full adder arranged across three carry-propagate stages. The registered input and output stages provide a fixed latency of 2 clock cycles and an ideal throughput of one multiplication result per clock cycle.

\subsection{Block Diagram}
\begin{figure}[H]
    \centering
    \includeimage{block_diagram.png}{4-bit Multiplier Structural Block Diagram}
    \caption{4-bit Array Multiplier Structural Architecture (16 AND Gates, 12 Full Adders, Input/Output Registers)}
    \label{fig:block_diagram}
\end{figure}

\subsection{Simulation and Verification}
Functional verification was carried out exhaustively using Icarus Verilog by cycling through all $2^{4} \times 2^{4} = 256$ input combinations ($a \in [0, 15], b \in [0, 15]$) with a self-checking testbench. A 2-cycle latency was accounted for before output sampling.

\begin{figure}[H]
    \centering
    \includeimage{partA_2.png}{Part A Simulation PASS Summary}
    \caption{Functional RTL Verification Summary (PASS on 256 vectors)}
    \label{fig:pass_result}
\end{figure}

\begin{figure}[H]
    \centering
    \includeimage{waveform.png}{Simulation Waveform with 2-Cycle Latency}
    \caption{Simulation Waveform showing Clock, Active-Low Reset, Inputs $a$, $b$, and Output Product $p$ with 2-Cycle Latency Marked}
    \label{fig:waveform}
\end{figure}

\section{Part B — Logic Synthesis with Yosys}

\subsection{Synthesis Results}
Logic synthesis was carried out using Yosys targeting the SkyWater 130~nm High-Density standard cell library (\texttt{sky130\_fd\_sc\_hd}). The design hierarchy was flattened and technology mapped using ABC and DFF library mapping. The area and cell count statistics extracted from \texttt{synth/mult\_array\_stat.rpt} are presented in Table~\ref{tab:partb_results}.

\begin{table}[H]
\centering
\begin{tabular}{|l|l|}
\hline
\textbf{Metric} & \textbf{Value} \\ \hline
Total number of cells & 68 \\ \hline
Number of flip-flops & 16 \\ \hline
Total cell area ($\mu\text{m}^2$) & 785.754 \\ \hline
Flip-flop area ($\mu\text{m}^2$) & 400.384 \\ \hline
Flip-flop area (\% of total) & 50.96\% \\ \hline
The 3 most-used cell types and counts & \begin{tabular}[c]{@{}l@{}}1. \texttt{sky130\_fd\_sc\_hd\_\_dfrtp\_1}: 16\\ 2. \texttt{sky130\_fd\_sc\_hd\_\_nand2\_1}: 8\\ 3. \texttt{sky130\_fd\_sc\_hd\_\_and2\_0}: 6 / \texttt{xnor2\_1}: 6\end{tabular} \\ \hline
Netlist simulation result & PASS: 256 vectors checked, 0 errors \\ \hline
\end{tabular}
\caption{Part B Logic Synthesis Results Summary (\texttt{mult\_array\_stat.rpt})}
\label{tab:partb_results}
\end{table}

\subsection{Gate-Level Netlist Simulation}
The synthesized gate-level Verilog netlist (\texttt{synth/mult\_array\_netlist.v}) was simulated against the functional testbench using the SKY130 standard cell primitives to ensure equivalence with RTL.

\begin{figure}[H]
    \centering
    \includeimage{netlist_pass.png}{Netlist Simulation PASS Screenshot}
    \caption{Gate-Level Netlist Simulation PASS Verification Output}
    \label{fig:netlist_pass}
\end{figure}

\section{Part C — Static Timing Analysis with OpenSTA}

\subsection{Timing Analysis Results}
Static timing analysis was conducted with OpenSTA using the typical process corner library (\texttt{sky130\_fd\_sc\_hd\_\_tt\_025C\_1v80.lib}) under an initial clock constraint of $T = 10.0$~ns with input/output delays set to $1.0$~ns. From the timing report (\texttt{sta/mult\_array\_timing.rpt}), the worst negative slack is $+7.027$~ns, giving a minimum achievable clock period of $T_{min} = 2.973$~ns ($\approx 2.97$~ns) and a maximum frequency of $F_{max} = 336.70$~MHz.

\begin{table}[H]
\centering
\begin{tabular}{|l|l|}
\hline
\textbf{Metric} & \textbf{Value} \\ \hline
Worst slack at $T = 10$ ns (ns) & $+7.03$ \\ \hline
$T_{min}$ (ns) & $2.97$ \\ \hline
$F_{max}$ (MHz) & $336.70$ \\ \hline
Slack at $T = T_{min}$ (ns) & $\approx 0.00$ ($-0.003$) \\ \hline
Critical path start point & \texttt{\_108\_} \\ \hline
Critical path end point & \texttt{\_103\_} \\ \hline
Number of cells on the critical path & 8 combinational gates (10 cells total) \\ \hline
Latency (ns) $= 2 \times T_{min}$ & $5.94$ \\ \hline
Throughput (million results/sec) & $336.70$ \\ \hline
\end{tabular}
\caption{Part C Static Timing Analysis Table}
\label{tab:partc_results}
\end{table}

\begin{figure}[H]
    \centering
    \includeimage{Part_C_Worst_Slack.png}{OpenSTA Worst Slack Output at T = 10 ns}
    \caption{OpenSTA Worst Slack Output at $T = 10$ ns}
    \label{fig:slack_10}
\end{figure}

\begin{figure}[H]
    \centering
    \includeimage{PartC_Corrected.png}{OpenSTA Worst Slack Output at T = T_min}
    \caption{OpenSTA Worst Slack Output at $T = T_{min}$}
    \label{fig:slack_zero}
\end{figure}

\subsection{Critical Path Cell Sequence}
The critical timing path originates at input register \texttt{\_108\_} and terminates at output register \texttt{\_103\_}. It propagates through the following 8 standard combinational logic cells:
\begin{enumerate}
    \item \texttt{\_108\_} (\texttt{sky130\_fd\_sc\_hd\_\_dfrtp\_1} - Start Flip-Flop, CLK $\rightarrow$ Q: $0.410$~ns)
    \item \texttt{\_050\_} (\texttt{sky130\_fd\_sc\_hd\_\_and4\_1}, Delay: $0.223$~ns)
    \item \texttt{\_061\_} (\texttt{sky130\_fd\_sc\_hd\_\_o21ai\_0}, Delay: $0.195$~ns)
    \item \texttt{\_062\_} (\texttt{sky130\_fd\_sc\_hd\_\_nand3\_1}, Delay: $0.140$~ns)
    \item \texttt{\_063\_} (\texttt{sky130\_fd\_sc\_hd\_\_maj3\_1}, Delay: $0.380$~ns)
    \item \texttt{\_069\_} (\texttt{sky130\_fd\_sc\_hd\_\_xnor3\_1}, Delay: $0.393$~ns)
    \item \texttt{\_074\_} (\texttt{sky130\_fd\_sc\_hd\_\_maj3\_1}, Delay: $0.369$~ns)
    \item \texttt{\_086\_} (\texttt{sky130\_fd\_sc\_hd\_\_maj3\_1}, Delay: $0.383$~ns)
    \item \texttt{\_087\_} (\texttt{sky130\_fd\_sc\_hd\_\_maj3\_1}, Delay: $0.349$~ns)
    \item \texttt{\_103\_} (\texttt{sky130\_fd\_sc\_hd\_\_dfrtp\_1} - End Flip-Flop, Setup: $-0.131$~ns)
\end{enumerate}
Total data arrival time through the path is $2.842$~ns. The 8 combinational cells precisely corroborate the theoretical critical path estimate of $2N = 8$ full-adder equivalent stages for an $N = 4$ array multiplier.

\section{Part D — Physical Design: RTL to GDSII}

\subsection{Implementation and Flow Results}
Physical design was completed using LibreLane on the SKY130 technology targeting a fixed die size of $150\,\mu\text{m} \times 150\,\mu\text{m}$ (\texttt{FP\_SIZING: absolute}). The automated flow executed floorplanning, power grid distribution (PDN), placement, clock tree synthesis (CTS), detailed routing, parasitic RC extraction, and signoff checks (DRC, LVS, STA). 

The metrics extracted from \texttt{pnr/runs/RUN\_2026-10-15\_20-59-19/final/metrics.json} are summarized in Table~\ref{tab:partd_results}.

\begin{table}[H]
\centering
\begin{tabular}{|l|l|}
\hline
\textbf{Metric} & \textbf{Value} \\ \hline
CLOCK\_PERIOD used (ns) & 10.0 \\ \hline
Number of runs needed & 1 \\ \hline
Core area ($\mu\text{m}^2$) & 17,759.50 (Die area: 22,500.00 $\mu\text{m}^2$) \\ \hline
Cells after synthesis (from Part B) & 68 \\ \hline
Cells after place and route & 1,892 (335 standard cells + 1,557 fill cells) \\ \hline
Kinds of cells added & \begin{tabular}[c]{@{}l@{}}Fill cells (1,557), Tap cells (245), \\ Timing buffers (19), Clock buffers (3), Inverters (4)\end{tabular} \\ \hline
Wirelength ($\mu\text{m}$) & 2,395 \\ \hline
Setup slack (ns) & +2.894 (\texttt{max\_ss} worst corner) / +6.475 (\texttt{nom\_tt}) \\ \hline
Post-route $F_{max}$ (MHz) & 140.73 (worst corner signoff) / 283.69 (\texttt{tt} corner) \\ \hline
$F_{max}$ from Part C (MHz) & 336.70 \\ \hline
DRC errors & 0 \\ \hline
LVS errors & 0 \\ \hline
\end{tabular}
\caption{Part D Physical Design Results Summary (\texttt{metrics.json})}
\label{tab:partd_results}
\end{table}

The post-route maximum frequency was computed according to the formula:
\begin{equation}
F_{max,\text{post-route}} = \frac{1000}{\text{CLOCK\_PERIOD} - \text{setup slack}} = \frac{1000}{10.0 - 2.894} = 140.73\text{ MHz}
\end{equation}
Under typical corner conditions (\texttt{nom\_tt\_025C\_1v80}), the setup slack is $+6.475$~ns, corresponding to an effective frequency of $283.69$~MHz.

\subsection{Layout Visualizations}
The final layout GDSII was inspected in KLayout to verify standard cell row arrangement, power stripe distribution, pin placement, and metal routing.

\begin{figure}[H]
    \centering
    \includeimage{klayout_whole_die.png}{KLayout Full Die View}
    \caption{Complete Die Layout in KLayout ($150\,\mu\text{m} \times 150\,\mu\text{m}$)}
    \label{fig:klayout_whole_die}
\end{figure}

\begin{figure}[H]
    \centering
    \includeimage{klayout_zoomed.png}{KLayout Zoomed View showing Cell Rows and Routing}
    \caption{Zoomed KLayout View showing Standard Cell Rows and Multi-layer Metal Interconnects}
    \label{fig:klayout_zoomed}
\end{figure}

\section{Questions for the Report}

\subsection{Question 1: Combinational Logic Cell Decomposition}
\textbf{Question:} The array needs 16 AND gates and 12 full adders. How many cells did synthesis produce for the logic (total cells minus flip-flops)? Explain why the number of cells is not simply $16 + 12$.

\textbf{Answer:} 
Synthesis produced exactly $68 - 16 = 52$ combinational logic cells. The number of logic cells is not simply $16 + 12 = 28$ because standard cell libraries such as SkyWater SKY130 do not contain single primitive macros for full adders. During logic optimization, Yosys flattens the structural full adder hierarchy, and ABC technology-maps the 1-bit full adder equations into elementary library gates—specifically 3-input majority gates (\texttt{maj3\_1}), 2-input and 3-input XOR/XNOR gates (\texttt{xnor2\_1}, \texttt{xnor3\_1}, \texttt{xor3\_1}), NAND gates, and complex AOI/OAI gates. A single full adder requires several standard cells to realize both its carry and sum outputs. Furthermore, ABC performs multi-level Boolean restructuring, logic factoring, and inverter insertion to minimize delay and area, resulting in 52 mapped logic cells rather than 28 idealized blocks.

\subsection{Question 2: Critical Path Identification and Theoretical Model}
\textbf{Question:} What is the critical path of your design, from which register to which register? How many cells are on it, and how does this compare with the estimate of about $2N$ full-adder delays?

\textbf{Answer:} 
The critical path begins at flip-flop \texttt{\_108\_} (an input operand register bit) and terminates at flip-flop \texttt{\_103\_} (the most significant product register bit), with a total data arrival time of $2.842$~ns. The path comprises exactly 8 combinational standard cells: \texttt{\_050\_} (\texttt{and4\_1}), \texttt{\_061\_} (\texttt{o21ai\_0}), \texttt{\_062\_} (\texttt{nand3\_1}), \texttt{\_063\_} (\texttt{maj3\_1}), \texttt{\_069\_} (\texttt{xnor3\_1}), \texttt{\_074\_} (\texttt{maj3\_1}), \texttt{\_086\_} (\texttt{maj3\_1}), and \texttt{\_087\_} (\texttt{maj3\_1}). In an $N \times N$ unsigned array multiplier, the worst-case propagation delay traverses down the diagonal partial-product generation and across the final ripple-carry stage, theoretically estimated as approximately $2N$ full-adder delays. For $N = 4$, $2N = 2 \times 4 = 8$ stages. The synthesized critical path consists of precisely 8 combinational gates, matching the theoretical estimate of $2N = 8$ with exact precision.

\subsection{Question 3: Sequential Area Dominance in Small Designs}
\textbf{Question:} What percentage of the total cell area is flip-flops? Why is this percentage large for such a small design?

\textbf{Answer:} 
Sequential elements account for $50.96\%$ of the total active cell area ($400.384\,\mu\text{m}^2$ out of $785.754\,\mu\text{m}^2$). This percentage is exceptionally large because standard edge-triggered D flip-flops (\texttt{sky130\_fd\_sc\_hd\_\_dfrtp\_1}) are physically large master-slave circuits comprising 20 to 26 transistors each, occupying an area of $25.024\,\mu\text{m}^2$ per cell. In contrast, simple combinational gates like 2-input NAND (\texttt{nand2\_1}) occupy only $3.754\,\mu\text{m}^2$. The architecture instantiates 16 flip-flops (4 for input $a$, 4 for input $b$, and 8 for product output $p$), consuming $16 \times 25.024\,\mu\text{m}^2 = 400.384\,\mu\text{m}^2$. Because the combinational datapath of a 4-bit multiplier is tiny ($52$ cells totaling only $385.37\,\mu\text{m}^2$), the bulky I/O pipeline registers dominate the design area.

\subsection{Question 4: Post-Route Frequency Degradation and Physical Cells}
\textbf{Question:} Why is the post-route $F_{max}$ lower than the $F_{max}$ from Part C? Which kinds of cells did the flow add, and why is each kind needed?

\textbf{Answer:} 
Post-route $F_{max}$ ($140.73$~MHz worst corner, $283.69$~MHz nominal) is lower than Part C pre-route $F_{max}$ ($336.70$~MHz) primarily due to physical interconnect parasitics and multi-corner signoff derating. While Part C assumes idealized zero-wire delays, post-routing includes distributed resistance and capacitance across $2,395\,\mu\text{m}$ of metal wiring, clock tree insertion delay and skew ($+0.265$~ns), and clock jitter/uncertainty ($0.250$~ns) under the worst-case slow-slow PVT corner (\texttt{nom\_ss\_100C\_1v60}). During physical implementation, the tool added four distinct categories of physical cells:
\begin{enumerate}
    \item \textbf{Tap cells} (245 cells, \texttt{tap\_cell}): Inserted at periodic pitch to tie n-wells to $V_{DD}$ and the substrate to $V_{SS}$, preventing CMOS latch-up.
    \item \textbf{Fill cells} (1,557 cells, \texttt{fill\_cell}): Inserted to fill vacant core sites, ensuring continuous power rails, well continuity, and metal density rules for chemical-mechanical polishing (CMP).
    \item \textbf{Timing repair and Clock buffers} (22 cells: 19 \texttt{timing\_repair\_buffer}, 3 \texttt{clock\_buffer}): Clock buffers balance the clock tree network across flip-flops to minimize clock skew; timing buffers drive long nets to fix maximum slew, transition, and setup/hold constraints.
    \item \textbf{Inverters} (4 cells, \texttt{inverter}): Provide polarity correction and load driving capability.
\end{enumerate}

\subsection{Question 5: DRC and LVS Signoff Verification}
\textbf{Question:} What does a DRC check verify, and what does an LVS check verify? Why must both pass before a chip can be manufactured?

\textbf{Answer:} 
Design Rule Checking (DRC) inspects the geometric features of all mask layers against the foundry\x27s manufacturing design rules (such as minimum metal widths, layer-to-layer spacings, enclosure margins, and antenna rules) to guarantee that the layout can be photolithographically manufactured without physical defects like opens or shorts. Layout Versus Schematic (LVS) extracts the electrical circuit (MOSFETs, diodes, capacitors, and their net connectivity) from physical layout polygons and performs a mathematical graph isomorphism comparison against the post-synthesis netlist to verify complete topological and schematic equivalence. Both checks must pass with zero violations because a DRC-clean layout may contain wiring errors that render the circuit non-functional (failing LVS), while an LVS-clean netlist with DRC violations will suffer catastrophic yield failure or physical short circuits during semiconductor fabrication.

\subsection{Question 6: Multiplier Scalability Analysis ($N = 8$)}
\textbf{Question:} Using the formulas in Section 2, estimate how the number of full adders and the critical path delay would change if the multiplier were 8 bits wide instead of 4. Show your working.

\textbf{Answer:} 
From Section 2 of the assignment guide, for an $N$-bit unsigned array multiplier:
\begin{equation}
\text{Full Adders: } FA(N) = N(N - 1)
\end{equation}
\begin{equation}
\text{AND Gates: } AND(N) = N^2
\end{equation}
\begin{equation}
\text{Critical Path Delay: } t_{crit}(N) \approx 2N \times t_{FA}
\end{equation}
For $N = 4$:
\begin{itemize}
    \item $FA(4) = 4(4 - 1) = 12$ full adders.
    \item $AND(4) = 4^2 = 16$ AND gates.
    \item Critical path delay corresponds to $2(4) = 8$ full-adder delays. With a measured combinational delay of $t_{comb} = 2.842 - 0.410 = 2.432$~ns, each full-adder equivalent stage contributes $\approx 0.304$~ns.
\end{itemize}
For $N = 8$:
\begin{itemize}
    \item $FA(8) = 8(8 - 1) = 56$ full adders. The number of full adders increases by a factor of $\frac{56}{12} \approx 4.67\times$ (an increase of 44 full adders).
    \item $AND(8) = 8^2 = 64$ AND gates, which is a $4\times$ increase.
    \item Critical path delay doubles to $2(8) = 16$ full-adder stages:
    \begin{equation}
    t_{comb, 8\text{-bit}} \approx 2 \times t_{comb, 4\text{-bit}} \approx 2 \times 2.432\text{ ns} = 4.864\text{ ns}
    \end{equation}
    Including register clock-to-Q ($0.410$~ns) and setup time ($-0.131$~ns), the estimated minimum clock period scales from $T_{min} = 2.973$~ns to:
    \begin{equation}
    T_{min, 8\text{-bit}} \approx 4.864 + 0.410 - 0.131 + \text{margin} \approx 5.40\text{ ns}
    \end{equation}
    This lowers the maximum pre-route clock frequency from $336.70$~MHz down to approximately $F_{max} \approx \frac{1000}{5.40} \approx 185.19$~MHz.
\end{itemize}

\end{document}
